\documentclass[11pt]{article}

\usepackage[utf8]{inputenc}
\usepackage[T1]{fontenc}
\usepackage{lmodern}
\usepackage{geometry}
\usepackage{amsmath,amssymb,amsfonts,amsthm,mathtools}
\usepackage{bm}
\usepackage{enumitem}
\usepackage{microtype}
\usepackage{hyperref}
\usepackage{titlesec}
\usepackage{booktabs}
\usepackage{array}
\usepackage{setspace}
\usepackage{mathrsfs}
\usepackage{physics}

\geometry{margin=1in}
\hypersetup{colorlinks=true, linkcolor=black, urlcolor=blue, citecolor=black}
\setlist[enumerate]{topsep=0.4em,itemsep=0.35em}
\setlist[itemize]{topsep=0.3em,itemsep=0.25em}
\titleformat{\section}{\large\bfseries}{\thesection.}{0.5em}{}
\titleformat{\subsection}{\normalsize\bfseries}{\thesubsection.}{0.5em}{}
\setstretch{1.06}

\newcommand{\Aether}{\AE{}ther}
\newcommand{\Aflow}{\AE{}ther-flow}
\newcommand{\TheoryName}{\Aflow{} Interpretation of Relativity}
\newcommand{\TheoryProgram}{\Aether{} / \Aflow{} framework}
\newcommand{\RespShapeReadout}{\widetilde q_{\mathrm{br}}}
\newcommand{\RespShapeCov}{\widetilde\kappa_{\mathrm{br}}}
\newcommand{\RespShapeRelax}{\widetilde\rho_{\mathrm{br}}}
\newcommand{\RespShapeWell}{\widetilde\lambda_{\mathrm{br}}}
\newcommand{\RespShapeEq}{\widetilde U_{\mathrm{br}}^{\ast}}
\newcommand{\RespShapeIso}{\widetilde\eta_{\mathrm{br}}}
\newcommand{\RespRawReadout}{\widehat q_{\mathrm{br}}}
\newcommand{\RespRawRef}{\widehat U_{\mathrm{br}}^{\mathrm{ref}}}
\newcommand{\RespRawCov}{\widehat\kappa_{\mathrm{br}}}
\newcommand{\RespRawRelax}{\widehat\rho_{\mathrm{br}}}
\newcommand{\RespRawWell}{\widehat\lambda_{\mathrm{br}}}
\newcommand{\RespRawEq}{\widehat U_{\mathrm{br}}^{\ast}}
\newcommand{\RespVacPhase}{\phi_{\mathrm{br}}}
\newcommand{\CurvProfileCan}{F_{\mathrm{br}}^{\mathrm{can}}}
\newcommand{\DownPkg}{\mathscr P_{\mathrm{down}}}
\newcommand{\ShapeTuple}{\mathfrak S_{\mathrm{br}}}
\newcommand{\RawTuple}{\mathfrak R_{\mathrm{br}}}
\newcommand{\ScaleMod}{s_{\mathrm{br}}}
\newcommand{\RespOrbitGrad}{\alpha_{\mathrm{br}}^{\mathrm{orb}}}
\newcommand{\RespOrbitEnergy}{\widetilde{\mathscr E}_{\mathrm{br}}^{\mathrm{orb}}}
\newcommand{\PrimClass}{\mathfrak C_{\mathrm{br}}}
\newcommand{\PrimRead}{r_{\mathrm{br}}}
\newcommand{\PrimCovSlot}{a_{\mathrm{br}}^{\varphi}}
\newcommand{\PrimRelaxSlot}{a_{\mathrm{br}}^{V}}
\newcommand{\PrimKinetic}{\bm a_{\mathrm{br}}}
\newcommand{\PrimWellSlot}{b_{\mathrm{br}}^{\lambda}}
\newcommand{\PrimEqSlot}{b_{\mathrm{br}}^{U}}
\newcommand{\PrimAnchor}{\bm b_{\mathrm{br}}}
\newcommand{\PrimMap}{\Pi_{\mathrm{shape}}}
\newcommand{\RawMap}{\Pi_{\mathrm{raw}}}
\newcommand{\DeepSector}{\mathfrak D_{\mathrm{br}}^{\mathrm{amp}}}
\newcommand{\DeepReadAmp}{\xi_{\mathrm{br}}^{q}}
\newcommand{\DeepCovAmp}{\xi_{\mathrm{br}}^{\varphi}}
\newcommand{\DeepRelaxAmp}{\xi_{\mathrm{br}}^{V}}
\newcommand{\DeepKinAmp}{\bm\xi_{\mathrm{br}}}
\newcommand{\DeepWellAmp}{\zeta_{\mathrm{br}}^{\lambda}}
\newcommand{\DeepEqAmp}{\zeta_{\mathrm{br}}^{U}}
\newcommand{\DeepAnchorAmp}{\bm\zeta_{\mathrm{br}}}
\newcommand{\DeepScaleAmp}{\varsigma_{\mathrm{br}}}
\newcommand{\AmpMap}{\Pi_{\mathrm{amp}}}
\newcommand{\DeepMap}{\Pi_{\mathrm{deep}}}
\newcommand{\SignQuot}{\sim_{\pm}}
\newcommand{\Rstar}{\mathbb R^{\times}}

\newtheorem{definition}{Definition}
\newtheorem{proposition}{Proposition}
\newtheorem{remark}{Remark}
\newtheorem{corollary}{Corollary}
\newtheorem{theorem}{Theorem}

\title{\textbf{The \TheoryName{}}\\[0.4em]
\large Nonlocal Bridge Self-Response Off-Shell Profile-Selection Bridge-Order Orbit-Shape/Modulus Primitive-Class Amplitude-Lift Origin Note}
\author{}
\date{}

\begin{document}

\maketitle

\begin{abstract}
The benchmark exact-theory statement of \TheoryName{} remains fixed and is not reopened here. The overview-first exact-closure package is still the public front door. The present note acts only on the optional deeper derivational bridge, and it uses the primitive-class setup note as a frozen interface.

That setup note already fixed the minimal positive coefficient class
\[
\PrimClass=(\PrimRead,\PrimKinetic,\PrimAnchor,\ScaleMod)
\in \mathbb R_{>0}\times\mathbb R_{>0}^2\times\mathbb R_{>0}^2\times\mathbb R_{>0},
\]
together with its map into the unit-reference orbit-shape block and then into the raw orbit. What it did not yet do was choose one concrete deeper sector below that class. The next honest step is therefore no longer another setup remark. It is the first direct still-deeper realization through that interface.

This note chooses the smallest such realization. The deeper sector is the signed-amplitude class
\[
\DeepSector
=
\Rstar\times(\Rstar)^2\times(\Rstar)^2\times\Rstar
\]
with coordinates
\[
(\DeepReadAmp,\DeepKinAmp,\DeepAnchorAmp,\DeepScaleAmp)
=
(\DeepReadAmp,(\DeepCovAmp,\DeepRelaxAmp),(\DeepWellAmp,\DeepEqAmp),\DeepScaleAmp).
\]
The map into the frozen interface is the componentwise square
\[
\AmpMap:
(\DeepReadAmp,\DeepKinAmp,\DeepAnchorAmp,\DeepScaleAmp)
\longmapsto
\bigl(
(\DeepReadAmp)^{2},
((\DeepCovAmp)^{2},(\DeepRelaxAmp)^{2}),
((\DeepWellAmp)^{2},(\DeepEqAmp)^{2}),
(\DeepScaleAmp)^{2}
\bigr).
\]
This realizes the primitive class as the quotient of \(\DeepSector\) by the finite independent sign-flip relation \(\SignQuot\).

Using the frozen interface exactly as fixed, the note then proves four facts. First, the amplitude lift reproduces the same unit-reference orbit-shape block and the same positive orbit modulus. Second, it reproduces the same raw-orbit lift, now written explicitly as squares of composite amplitudes. Third, it reproduces the same downstream density/current block, the same composite primitive bridge-order block, the same phase-neutral minimizing branch, the same canonical profile, and the same dressed bridge map. Fourth, the new quotient structure is only the finite sign quotient of the square-root amplitudes themselves. It neither identifies distinct positive moduli nor locks the global \(O(2)\) vacuum angle. Therefore the current verdicts remain unchanged: the positive scale orbit is still a canonical-normalization orbit rather than a proved redundancy, and the vacuum angle remains residual.

The scope is deliberately narrow. This is the first direct deeper realization below the primitive-class interface, not a claim that the final unique substrate microphysics has already been reached.
\end{abstract}

\section{Introduction}

The active repository no longer needs another bridge note in order to justify its public theory claim. The exact-closure sequence already presents \TheoryName{} as the disciplined exact relativistic theory statement built on the \Aether{} / \Aflow{} ontology, and that benchmark package remains the front-door reading order. The present note therefore does not strengthen the flagship theory claim. It acts only on the optional deeper derivational continuation beneath that fixed benchmark.

On that optional continuation line, two recent notes already fixed the discipline of the next move. The orbit-shape/modulus-origin note fixed the burden: derive or sharply constrain the unit-reference orbit-shape block and the positive orbit modulus while preserving the same downstream chain. The primitive-class setup note then fixed the minimal interface below that burden: every later deeper derivation must factor through the readout block, the ordered kinetic pair, the radial anchor pair, and the positive orbit modulus.

That setup note was necessary, but it was not yet the first direct deeper manuscript. The still-open choice was the concrete deeper sector feeding that interface. The present note makes that choice explicitly and minimally.

\paragraph{Framework and claim boundary.}
\TheoryProgram{} denotes the broader \Aether{} / \Aflow{} framework built on the claims that the \Aether{} is the underlying four-dimensional substrate of reality, the \Aflow{} is its intrinsic ordered motion, observed three-dimensional space is the local experiential slice of that deeper substrate, S-time is the experienced order of change arising from matter, light, and the \Aflow{}, and observed expansion is the three-dimensional appearance of deeper four-dimensional ordered motion. Within that framework, \TheoryName{} denotes the exact-closure theory in which the effective gravitational dynamics are adopted to be exactly Einsteinian with universal matter coupling. In the active sequence, \emph{adoption} means use of that established relativistic dynamics without claiming substrate derivation, while \emph{derivation} is reserved for a first-principles recovery from explicit substrate variables.

This note stays strictly inside that boundary. Its positive claim is not that a unique final substrate completion has already been isolated. Its positive claim is narrower: the first direct deeper realization below the frozen primitive-class interface can now be written explicitly, can be checked coefficient by coefficient, and can be shown to preserve the same already-positive downstream package without changing the current scale-orbit or vacuum-angle verdicts.

\section{Frozen Interface and Concrete Deeper Sector}

\begin{definition}[Frozen primitive-class interface]
\label{def:interface}
At the interface fixed by the setup note, the primitive class is
\begin{equation}
\PrimClass
:=
(\PrimRead,\PrimKinetic,\PrimAnchor,\ScaleMod)
\in
\mathbb R_{>0}\times\mathbb R_{>0}^{2}\times\mathbb R_{>0}^{2}\times\mathbb R_{>0},
\label{eq:primclass}
\end{equation}
with
\begin{equation}
\PrimKinetic=(\PrimCovSlot,\PrimRelaxSlot),
\qquad
\PrimAnchor=(\PrimWellSlot,\PrimEqSlot),
\label{eq:blocks}
\end{equation}
and the interface map
\begin{equation}
\PrimMap(\PrimRead,\PrimKinetic,\PrimAnchor,\ScaleMod)
=
(\RespShapeReadout,\RespShapeCov,\RespShapeRelax,\RespShapeWell,\RespShapeEq,\ScaleMod)
\label{eq:primmap}
\end{equation}
given by
\begin{equation}
\RespShapeReadout=\PrimRead,
\qquad
\RespShapeCov=\PrimCovSlot,
\qquad
\RespShapeRelax=\PrimRelaxSlot,
\label{eq:shapeA}
\end{equation}
\begin{equation}
\RespShapeWell=\PrimWellSlot,
\qquad
\RespShapeEq=\PrimEqSlot.
\label{eq:shapeB}
\end{equation}
The raw-orbit lift is then
\begin{equation}
\RespRawReadout=\ScaleMod\,\RespShapeReadout,
\qquad
\RespRawRef=\ScaleMod,
\qquad
\RespRawCov=\ScaleMod^{2}\RespShapeCov,
\label{eq:rawliftA}
\end{equation}
\begin{equation}
\RespRawRelax=\ScaleMod^{2}\RespShapeRelax,
\qquad
\RespRawWell=\RespShapeWell,
\qquad
\RespRawEq=\ScaleMod\,\RespShapeEq.
\label{eq:rawliftB}
\end{equation}
\end{definition}

\begin{definition}[Amplitude-lift deeper sector]
\label{def:deep}
Write \(\Rstar:=\mathbb R\setminus\{0\}\). The concrete deeper sector chosen in this note is
\begin{equation}
\DeepSector
:=
\Rstar\times(\Rstar)^{2}\times(\Rstar)^{2}\times\Rstar,
\label{eq:deepsector}
\end{equation}
with coordinates
\begin{equation}
(\DeepReadAmp,\DeepKinAmp,\DeepAnchorAmp,\DeepScaleAmp)
=
(\DeepReadAmp,(\DeepCovAmp,\DeepRelaxAmp),(\DeepWellAmp,\DeepEqAmp),\DeepScaleAmp).
\label{eq:deepcoords}
\end{equation}
\end{definition}

\begin{remark}
Definition \ref{def:deep} is the smallest direct deeper choice compatible with the frozen interface: each positive interface slot is fed by one signed square-root amplitude. The note does not claim that these amplitudes are already the final microscopic variables of the \Aether{}. It claims only that they are a concrete and disciplined first deeper sector below the setup layer.
\end{remark}

\section{Amplitude Lift Into the Primitive Class}

\begin{definition}[Amplitude-lift map]
\label{def:ampmap}
Define
\begin{equation}
\AmpMap:\DeepSector\longrightarrow\PrimClass
\label{eq:ampmap}
\end{equation}
by
\begin{equation}
\AmpMap(\DeepReadAmp,\DeepKinAmp,\DeepAnchorAmp,\DeepScaleAmp)
=
\bigl(
(\DeepReadAmp)^{2},
((\DeepCovAmp)^{2},(\DeepRelaxAmp)^{2}),
((\DeepWellAmp)^{2},(\DeepEqAmp)^{2}),
(\DeepScaleAmp)^{2}
\bigr).
\label{eq:ampmapexplicit}
\end{equation}
\end{definition}

\begin{definition}[Independent sign quotient]
\label{def:signquot}
Define an equivalence relation \(\SignQuot\) on \(\DeepSector\) by
\begin{align}
&(\DeepReadAmp,\DeepCovAmp,\DeepRelaxAmp,\DeepWellAmp,\DeepEqAmp,\DeepScaleAmp)
\SignQuot
(\epsilon_q\DeepReadAmp,\epsilon_\varphi\DeepCovAmp,\epsilon_V\DeepRelaxAmp,
\epsilon_\lambda\DeepWellAmp,\epsilon_U\DeepEqAmp,\epsilon_s\DeepScaleAmp)
\label{eq:signquot}
\end{align}
for arbitrary independent signs
\(\epsilon_q,\epsilon_\varphi,\epsilon_V,\epsilon_\lambda,\epsilon_U,\epsilon_s\in\{\pm1\}\).
\end{definition}

\begin{proposition}[The amplitude map is constant on sign orbits]
\label{prop:constant}
If two points of \(\DeepSector\) are related by \(\SignQuot\), then they have the same image under \(\AmpMap\).
\end{proposition}

\begin{proof}
Every component of \(\AmpMap\) is a square. Independent sign changes therefore leave all six squared components unchanged.
\end{proof}

\begin{theorem}[Primitive class as amplitude-lift quotient]
\label{thm:quotient}
The amplitude-lift map \(\AmpMap\) descends to a bijection
\begin{equation}
\DeepSector/\SignQuot \;\cong\; \PrimClass.
\label{eq:quotient}
\end{equation}
Equivalently, every point of \(\PrimClass\) has a unique amplitude lift modulo independent sign flips.
\end{theorem}

\begin{proof}
Surjectivity is immediate: given any
\(
(\PrimRead,(\PrimCovSlot,\PrimRelaxSlot),(\PrimWellSlot,\PrimEqSlot),\ScaleMod)\in\PrimClass,
\)
choose the positive square-root representative
\[
(\sqrt{\PrimRead},(\sqrt{\PrimCovSlot},\sqrt{\PrimRelaxSlot}),(\sqrt{\PrimWellSlot},\sqrt{\PrimEqSlot}),\sqrt{\ScaleMod})
\in\DeepSector.
\]
Its image under \(\AmpMap\) is exactly the given primitive tuple.

For injectivity modulo \(\SignQuot\), suppose two amplitude tuples have the same image under \(\AmpMap\). Then the six corresponding squared components agree. Over \(\Rstar\), equality of squares implies equality up to sign in each component separately. Hence the two amplitude tuples differ by independent sign flips and therefore lie in the same \(\SignQuot\)-class. This proves \eqref{eq:quotient}.
\end{proof}

\begin{corollary}[The first direct deeper note is now explicit]
\label{cor:explicit}
The present amplitude-lift sector is a concrete deeper sector below the setup note's frozen interface. It realizes the primitive class exactly, with the only redundancy being the finite sign quotient \(\SignQuot\).
\end{corollary}

\begin{proof}
This is the content of Theorem \ref{thm:quotient}.
\end{proof}

\section{Map Into Orbit-Shape/Modulus Data and the Raw Lift}

\begin{definition}[Deeper map into the orbit-shape/modulus data]
\label{def:deepmap}
Define
\begin{equation}
\DeepMap:=\PrimMap\circ\AmpMap:\DeepSector\longrightarrow(\ShapeTuple,\ScaleMod).
\label{eq:deepmap}
\end{equation}
Explicitly,
\begin{equation}
\RespShapeReadout=(\DeepReadAmp)^{2},
\qquad
\RespShapeCov=(\DeepCovAmp)^{2},
\qquad
\RespShapeRelax=(\DeepRelaxAmp)^{2},
\label{eq:deepshapeA}
\end{equation}
\begin{equation}
\RespShapeWell=(\DeepWellAmp)^{2},
\qquad
\RespShapeEq=(\DeepEqAmp)^{2},
\qquad
\ScaleMod=(\DeepScaleAmp)^{2}.
\label{eq:deepshapeB}
\end{equation}
\end{definition}

\begin{theorem}[Same raw-orbit lift from the amplitude sector]
\label{thm:rawlift}
For every point of \(\DeepSector\), the induced raw tuple is exactly the frozen-interface lift \eqref{eq:rawliftA}--\eqref{eq:rawliftB}. In amplitude variables this becomes
\begin{equation}
\RespRawReadout=(\DeepScaleAmp\DeepReadAmp)^{2},
\qquad
\RespRawRef=(\DeepScaleAmp)^{2},
\qquad
\RespRawCov=(\DeepScaleAmp\DeepCovAmp)^{2},
\label{eq:rawampA}
\end{equation}
\begin{equation}
\RespRawRelax=(\DeepScaleAmp\DeepRelaxAmp)^{2},
\qquad
\RespRawWell=(\DeepWellAmp)^{2},
\qquad
\RespRawEq=(\DeepScaleAmp\DeepEqAmp)^{2}.
\label{eq:rawampB}
\end{equation}
\end{theorem}

\begin{proof}
Substitute \eqref{eq:deepshapeA}--\eqref{eq:deepshapeB} into the frozen-interface formulas \eqref{eq:rawliftA}--\eqref{eq:rawliftB}. The result is exactly \eqref{eq:rawampA}--\eqref{eq:rawampB}.
\end{proof}

\begin{remark}
Theorem \ref{thm:rawlift} is the first direct deeper realization asked for by the setup layer. The raw orbit is not changed, repaired, or reinterpreted. It is reproduced exactly, now as the image of one explicit deeper signed-amplitude sector.
\end{remark}

\section{Same Phase-Neutral Branch and Same Downstream Package}

\begin{definition}[Induced orbit coefficients]
\label{def:iso}
For any amplitude datum in \(\DeepSector\), define the induced orbit coefficient
\begin{equation}
\RespShapeIso
:=
\frac{\RespShapeCov\,\RespShapeRelax}{\RespShapeCov+\RespShapeRelax}
=
\frac{(\DeepCovAmp)^{2}(\DeepRelaxAmp)^{2}}{(\DeepCovAmp)^{2}+(\DeepRelaxAmp)^{2}}.
\label{eq:iso}
\end{equation}
\end{definition}

\begin{theorem}[Same phase-neutral minimizing branch]
\label{thm:branch}
The orbit-level energy induced by Definition \ref{def:deepmap} is the same frozen-interface energy as in the setup and raw-orbit notes, now with coefficients \eqref{eq:deepshapeA}--\eqref{eq:deepshapeB}. Its distinguished minimizing constant branch is therefore
\begin{equation}
\ScaleMod(x)\equiv (\DeepScaleAmp)^{2},
\qquad
U(x)\equiv (\DeepScaleAmp\DeepEqAmp)^{2},
\qquad
V(x)\equiv 0,
\qquad
\RespVacPhase(x)\equiv \phi_{0},
\label{eq:branch}
\end{equation}
for arbitrary constant \(\phi_{0}\in\mathbb R/2\pi\mathbb Z\). In particular, the minimizing branch remains phase-neutral.
\end{theorem}

\begin{proof}
By Definition \ref{def:deepmap}, the amplitude sector produces exactly the same orbit-shape coefficients and positive modulus used by the frozen-interface energy. The minimizing-family computation in the raw-vacuum-orbit layer depends only on those induced coefficients and on the lift formulas \eqref{eq:rawliftA}--\eqref{eq:rawliftB}. By Theorem \ref{thm:rawlift}, those formulas are reproduced exactly. Therefore the minimizing constant family is unchanged, and in amplitude variables it is written precisely as \eqref{eq:branch}. Because \(\phi_{0}\) remains arbitrary and constant while \(V\equiv0\), the branch is still phase-neutral.
\end{proof}

\begin{definition}[Frozen downstream package]
\label{def:downstream}
At the present derivational stage, the downstream package \(\DownPkg\) means the same already-recorded density/current block, the same composite primitive bridge-order block, the same phase-neutral minimizing branch, the same canonical profile, and the same dressed bridge map, all understood to depend on deeper data only through the induced pair \((\ShapeTuple,\ScaleMod)\) and the same minimizing branch.
\end{definition}

\begin{theorem}[Same downstream density/current block, bridge-order block, canonical profile, and dressed bridge map]
\label{thm:down}
For the amplitude-lift deeper sector \(\DeepSector\), the downstream package \(\DownPkg\) is recovered unchanged. In particular:
\begin{enumerate}
    \item the same downstream density/current block is recovered;
    \item the same composite primitive bridge-order block is recovered;
    \item the same phase-neutral minimizing branch of Theorem \ref{thm:branch} is recovered;
    \item the same canonical profile is recovered, now written as
    \begin{equation}
    \CurvProfileCan(x)
    =
    1+\RespShapeEq\,x^{2}
    =
    1+(\DeepEqAmp)^{2}x^{2}
    =
    1+\frac{\RespRawEq}{\RespRawRef}x^{2};
    \label{eq:profile}
    \end{equation}
    \item the same dressed bridge map is recovered.
\end{enumerate}
\end{theorem}

\begin{proof}
Definition \ref{def:deepmap} shows that the amplitude sector feeds the frozen interface only through the induced orbit-shape/modulus data \((\ShapeTuple,\ScaleMod)\). Theorem \ref{thm:rawlift} shows that the same raw tuple is then generated. Theorem \ref{thm:branch} shows that the same phase-neutral minimizing branch survives. By Definition \ref{def:downstream}, all recorded downstream constructions depend on the deeper layer only through exactly those data. Therefore every downstream object in \(\DownPkg\) is unchanged.

The canonical profile formula \eqref{eq:profile} is immediate from \(\RespShapeEq=(\DeepEqAmp)^{2}\) and \(\RespRawEq/\RespRawRef=\RespShapeEq\).
\end{proof}

\begin{remark}
Theorem \ref{thm:down} is why the amplitude-lift note is honest at current scope. It deepens the origin of the interface without trying to earn novelty by rewriting downstream formulas that are already fixed positively.
\end{remark}

\section{What the New Quotient Does and Does Not Change}

\begin{definition}[Relevant verdict-changing structure]
\label{def:relevant}
At the present stage, a deeper sector changes the current scale-orbit or residual vacuum-angle verdicts only if it provides one of the following:
\begin{enumerate}
    \item a quotient identifying distinct positive moduli \(\ScaleMod^{(1)}\neq \ScaleMod^{(2)}\);
    \item a locking or quotient mechanism fixing or identifying distinct constant vacuum angles \(\RespVacPhase=\phi_{0}\).
\end{enumerate}
\end{definition}

\begin{proposition}[The amplitude-lift quotient is internal to square-root choice]
\label{prop:internal}
The only quotient carried by \(\DeepSector\) is the finite sign quotient \(\SignQuot\) of Definition \ref{def:signquot}. It identifies different signed square-root representatives of the same positive interface data, but it does not identify distinct positive values of \(\ScaleMod\).
\end{proposition}

\begin{proof}
By construction, \(\SignQuot\) changes only signs of the six amplitude coordinates while leaving their squares fixed. In particular, it relates \(\DeepScaleAmp\) only to \(-\DeepScaleAmp\), both of which determine the same positive modulus \((\DeepScaleAmp)^{2}\). If \((\DeepScaleAmp^{(1)})^{2}\neq(\DeepScaleAmp^{(2)})^{2}\), then no sign choice can relate the two amplitudes. Hence the quotient is internal to square-root representation and does not identify distinct positive moduli.
\end{proof}

\begin{theorem}[Scale-orbit and residual vacuum-angle verdicts remain unchanged]
\label{thm:verdict}
The amplitude-lift deeper sector introduces no verdict-changing structure in the sense of Definition \ref{def:relevant}. Therefore:
\begin{enumerate}
    \item the positive scale orbit remains a canonical-normalization orbit rather than a proved redundancy;
    \item the global \(O(2)\) vacuum angle remains residual.
\end{enumerate}
\end{theorem}

\begin{proof}
By Proposition \ref{prop:internal}, the only new quotient is the finite sign quotient on square-root representatives, and it does not identify distinct positive moduli. Thus no positive-modulus quotient has been added, so item (1) follows.

Likewise, the amplitude sector \(\DeepSector\) introduces no new angular variable and no anisotropic locking term. The minimizing family of Theorem \ref{thm:branch} still carries an arbitrary constant \(\phi_{0}\), and the amplitude-lift map does not identify different values of \(\phi_{0}\). Therefore the global \(O(2)\) vacuum angle remains residual, giving item (2).
\end{proof}

\begin{remark}
Theorem \ref{thm:verdict} is intentionally strict. The present deeper note \emph{does} introduce a real quotient, but it is only the finite sign quotient of square-root amplitudes. That quotient is too small and too internal to alter the recorded positive scale-orbit or vacuum-angle verdicts.
\end{remark}

\section{Claim Boundary}

This note does not claim that the final substrate microphysics below the primitive class has already been found uniquely. It does not claim that the signed amplitude variables of Definition \ref{def:deep} are the only possible deeper realization, and it does not claim that the positive scale orbit or the global \(O(2)\) vacuum angle have already been removed by a deeper redundancy or locking mechanism.

Its claim is narrower. Using the primitive-class setup note as a frozen interface, the present manuscript chooses the first concrete deeper sector below that interface, gives the explicit map into the readout block, ordered kinetic pair, radial anchor pair, and positive orbit modulus, and proves that the same raw-orbit lift and the same downstream package are recovered exactly. The only new quotient is the finite sign quotient of square-root amplitudes themselves, and that quotient does not change the current verdicts on the positive scale orbit or the residual vacuum angle.

The benchmark exact-theory package remains untouched. The already-positive downstream bridge package remains untouched. What is advanced here is only the first direct deeper realization beneath the setup layer.

\section{Conclusion}

The next honest manuscript step below the primitive-class setup note was to stop speaking only at interface level and to choose one explicit deeper sector. This note now does that. The chosen sector is the minimal signed amplitude lift: one signed square-root amplitude for the readout block, two for the ordered kinetic pair, two for the radial anchor pair, and one for the positive orbit modulus.

The positive result is exact and limited. Modulo the finite independent sign quotient, that amplitude sector reproduces the primitive class fixed by the setup note. Through the frozen interface it then reproduces the same unit-reference orbit-shape block, the same raw-orbit lift, the same downstream density/current block, the same primitive bridge-order block, the same phase-neutral minimizing branch, the same canonical profile, and the same dressed bridge map.

The scope is equally clear. The new quotient is only the sign quotient of square-root representatives. It does not identify different positive moduli and it does not lock or quotient the global \(O(2)\) vacuum angle. Therefore the disciplined current verdict remains in force: the positive scale orbit is still a canonical-normalization orbit rather than a proved deeper redundancy, and the vacuum angle remains residual. The next honest open burden, if the derivational line continues, is to derive or motivate the amplitude variables themselves from a still deeper substrate sector rather than leaving the amplitude lift as the present deepest recorded layer.

\end{document}
